\chapter*{How to read this foundation}
\addcontentsline{toc}{chapter}{Preface}
\label{ch:preface}


Construct the objects, prove their laws, and compare different computations of the same value.


This book begins with familiar questions: how do we compute \(\sqrt2\) or \(\pi\), and why do those computations work? They lead naturally to functions, integrals, power series, and exact identities. The aim is a foundation for mathematics in which the computations and their justifications stay visible. Decimal evaluations are available when requested; the main results explain what the constructed objects are and how they behave.



\section*{A number comes with a computation}
\label{preface:objects}


A represented real starts with a sequence of rational endpoint pairs:

\[n\longmapsto I_n=[a_n,b_n],\qquad a_n,b_n\in\mathbb Q.\]

It becomes valid only after we prove that the endpoints are ordered, later intervals lie inside earlier ones, and their widths tend to zero. These are separate mathematical obligations. No prescribed rate belongs to the definition; a concrete construction can supply a rate or a procedure for reaching a requested tolerance.


The field \texttt{RealRaw.compute} is a mathematical function. Its type alone does not certify that it is executable. Our concrete constructions use finite rational algorithms, and their operations must preserve executability for executable inputs. General laws may nevertheless quantify over arbitrary valid represented inputs. “Computable” describes the construction discipline; it is not a restriction that silently removes irrational coefficients or inputs from an exact theorem.


The foundation uses its own represented reals and complexes. It does not import Mathlib’s real or complex numbers, or use an abstract completion to supply missing objects.



\section*{An exact identity compares computations}
\label{preface:equality}


Two valid representations can have different intervals and different rates while representing the same value. We write

\[x\simeq y\]

for \texttt{RealRaw.Equiv}: their intervals overlap at every common stage. Validity matters here. Nesting and shrinking make this overlap relation transitive on valid representations. Refinement schedules are proof tools, rather than an extra part of the definition of equality.


This is an exact statement about every stage, not a claim that two long decimal strings happen to agree. Immediately after the first square-root construction, for example, we prove

\[s\cdot s\simeq2.\]

The squared rational bounds contain \(2\) at every stage. This identifies the computed number mathematically, before any request for digits. A decimal approximation still has its usual precision convention: rounding to \(k\) decimal places has error at most \(\tfrac12 10^{-k}\), while truncation determines an interval of width \(10^{-k}\).



\section*{Rational-input functions are primary}
\label{preface:functions}


A square-root computation with rational parameter \(q\) constructs the function \(q\mapsto\sqrt q\); \(\sqrt2\) is one evaluation. Likewise, \(\pi=4\arctan(1)\) leads us from a constant to arctangent. The primary function takes rational inputs, rational complex inputs, or rational coordinates in an explicitly described chart. These inputs are its domain, rather than samples of a function already supplied on a completed continuum.


Evaluation at a represented irrational input is a derived add-on. Continuity gives the bridge: for each positive rational output tolerance, provide an input tolerance controlling the change in the function, and account for the evaluation errors. Shrinking input intervals then justify the output construction. Agreement under equivalent input representations also needs proof.


Different computations can use different coordinates. A rational multiple of \(\pi\), a rational radian, and a rational circle parameter are different inputs. We compare on an overlap, along a justified coordinate transformation, or after constructing extensions. Agreement on overlapping domains does not by itself give a globally transitive relation between functions with different domains. There is no need to discard one computation in favor of a single privileged definition.



\section*{Derive with a skill; certify with a proof}
\label{preface:derive-check}


This separation takes some getting used to. A skill can use substitution, partial fractions, or a familiar trigonometric manipulation to propose the integral of a particular function. We need not first formalize that entire technique as a general Lean theorem. Reusable checked lemmas are helpful, but the derivation can remain a guide for discovering the answer.


The proposed answer then earns its mathematical meaning through a concrete proof. For a definite integral, construct the endpoint expression, verify its derivative against the independently specified integrand on the whole segment, and use a justified FTC argument:

\[\int_a^x f(t)\,dt\simeq F(x)-F(a).\]

The derivative check is often easier than formalizing the technique that found \(F\). The subtraction fixes the base point. Pole avoidance, branch conditions, and validity must hold throughout the segment; checking only the endpoints cannot justify crossing a singularity. The derivative identity becomes an integral identity through FTC and the required integral evidence.


For example, our concrete square-root secant proof certifies

\[\int_a^b\frac{dt}{2\sqrt t}\simeq\sqrt b-\sqrt a\qquad(0<a\le b).\]

Skills guide discovery; Lean checks the proposed result. This is a procedure for proof development, not a claim that every function has a closed form or that an automatic integrator is already implemented. \href{https://liuyao12.github.io/computable-analysis/skill-real-integrals.html}{Read the integration skill}, or \href{https://liuyao12.github.io/computable-analysis/foundation-story/index.html\#ftc}{inspect the checked square-root route}.



\section*{General laws and particular constructions}
\label{preface:laws}


A law can start with supplied, justified objects and prove a new consequence. Proving that such objects exist for a particular function is a separate job. Neither a record named “integral” nor an assumed endpoint identity supplies that proof. Each function brings its own construction and evidence; we do not first insist that every integrand belong to one preselected class.


The current \texttt{Integral.HasIntegral} API uses arbitrarily tight rectangle bounds. This is a checked sufficient route, more restrictive than the intended method-independent core. Numerical validity alone does not establish an integral’s role. The preface describes the intended organization and identifies this implementation boundary explicitly.


The same discipline applies to inverses and derivatives. Convexity or concavity can suggest a useful interval update without being formalized as a general prerequisite. For square root, explicit supporting-line and secant formulas suggest an algorithm; rational inequalities prove containment, nesting, and shrinking. For the functions we actually construct, neighboring secants and their error bounds construct the derivatives we need.



\section*{What “computable” means here}
\label{preface:terminology}


There is substantial common ground with established computable analysis. In the represented-spaces approach, points have names, and computability of a map is defined through a computable transformation of those names. This framework also accommodates noncomputable points. Rational interval representations are therefore not, by themselves, an unfamiliar idea. \href{https://arxiv.org/html/1204.3763\#S3}{Pauly’s introduction} explains that framework.


The distinctive choices of this project are the primary rational-input functions, the separate extension proofs, exact comparison of retained computations, and the use of function-specific constructions beneath general laws. We are developing analysis with explicit computations as its foundation; we are not currently presenting a Turing-machine classification of which analytic problems are computable.


Bishop-style constructive mathematics places a restriction on proof reasoning, typically expressed through intuitionistic logic. Our repository permits classical reasoning: the checked theorem audits may include \texttt{Classical.choice}. We separately require the actual constructors to be executable and their intervals to be justified. Thus a classical proof of validity does not automatically extract an algorithm, and classical choice must not substitute for a missing executable construction. \href{https://researchcommons.waikato.ac.nz/items/5e0e94b5-e729-445a-b11d-56a968dac2cd}{Bridges and Reeves} discuss the constructive tradition.


These distinctions concern the objects, interfaces, and proof discipline. They do not claim that interval arithmetic is new, or that constructive analysis cannot develop completeness. A connection to an established framework would itself be a theorem to prove.



\section*{The route through the opening chapters}
\label{preface:reading}


Chapter 1 begins with numbers and discovers rational-input functions. Geometry supplies circle and angle computations. Arctangent then motivates integration and power series. Concrete secant estimates introduce derivatives when we need FTC and Taylor arguments to compare computations. Specific inversion of arctangent constructs sine and cosine without a general inverse function theorem.


The \href{https://liuyao12.github.io/computable-analysis/foundation-story/index.html}{expandable construction map} separates this mathematical route from the actual declaration dependencies extracted from Lean. Checked statements, quantitative intermediate results, and missing comparison bridges are identified at their stated scope. A narrative explanation does not claim a formal theorem that has not been proved.


Everyone is welcome to join: mathematical ideas, Lean proofs, explanations, examples, and improvements to the book all help. You can work with or without an AI agent. \href{https://github.com/liuyao12/computable-analysis\#join-the-collaboration}{Join the collaboration}.


\href{https://liuyao12.github.io/computable-analysis/ch-foundations.html\#numbers-to-functions}{Begin Chapter 1: numbers lead to functions}

